\documentclass[../../main.tex]{subfiles}
\begin{document}

\begin{flushright}
  \footnotesize\textit{【自動文字起こし・要確認】}
\end{flushright}

\noindent
{\bf [解答]} ベクトルの2つの公理を①, ②とする. すなわち,
\begin{align}
  \vec{A} + \vec{B} = \vec{B} + \vec{A} \quad \cdots \text{①}
\end{align}
\begin{align}
  (\vec{A} + \vec{B}) + \vec{C} = \vec{A} + (\vec{B} + \vec{C}) \quad \cdots \text{②}
\end{align}
以下, $\vec{S}_n = \sum_{k=1}^n \vec{A}_k$, $\vec{T}_n = \sum_{k=1}^n \vec{B}_k$ とする. $\vec{S}_n = \vec{T}_n \quad \cdots$ ③ を帰納的に示す.

\paragraph{[補題1]}
任意のベクトル $\vec{a}_1, \dots, \vec{a}_n$ に対し ($n \ge 2$)
\begin{align}
  \vec{a}_1 + \dots + \vec{a}_n = \vec{a}_2 + \vec{a}_3 + \dots + \vec{a}_n + \vec{a}_1 \quad \cdots \star
\end{align}

(証) 帰納法による. $n=2$ の時は $\star$ は ①に他ならない. $n=k$ での成立を仮定すると, $n=k+1$ の時
\begin{align}
  \vec{a}_1 + \dots + \vec{a}_{k+1} &= \vec{a}_1 + \dots + \vec{a}_{k-1} + (\vec{a}_k + \vec{a}_{k+1}) \quad (\because \text{②}) \\
  &= \vec{a}_2 + \dots + (\vec{a}_k + \vec{a}_{k+1}) + \vec{a}_1 \quad (\because \text{仮定}) \\
  &= \vec{a}_2 + \dots + \vec{a}_k + \vec{a}_{k+1} + \vec{a}_1 \quad (\because \text{②})
\end{align}
ゆえに $n=k+1$ でも $\star$ は成立.
以上から $\star$ は示された. \hfill 終

\bigskip
まず, $n=1$ の時. ③ $\iff \vec{A}_1 = \vec{B}_1$ でこれは定義から成立する. そこで $i \in \mathbb{N}$ に対し, $n \le i$ での③の成立を仮定する. $n=i+1$ の時, $k=0, 1, \dots, i$ に対して,
\begin{align}
  \vec{A}_{i+1} = \vec{B}_{k+1}
\end{align}
なる $k$ があることに注意する.

\begin{enumerate}
  \item[$1^\circ$] $k=i$ の時
    \begin{align}
      \vec{T}_{i+1} &= \vec{T}_i + \vec{B}_{i+1} \\
      &= \vec{S}_i + \vec{A}_{i+1} \quad (\because \text{仮定}, \vec{A}_{i+1} = \vec{B}_{i+1}) \\
      &= \vec{S}_{i+1}
    \end{align}
    となり, $n=i+1$ でも ③は成立.

  \item[$2^\circ$] $k \ne i$ の時
    \begin{align}
      \vec{T}_{i+1} &= \vec{B}_1 + \dots + \vec{B}_k + \vec{B}_{k+1} + \vec{B}_{k+2} + \dots + \vec{B}_{i+1} \\
      &= \vec{B}_{k+1} + \vec{B}_1 + \dots + \vec{B}_k + \vec{B}_{k+2} + \dots + \vec{B}_{i+1} \quad (\because \text{仮定}) \\
      &= \vec{B}_1 + \dots + \vec{B}_k + \vec{B}_{k+2} + \dots + \vec{B}_{i+1} + \vec{B}_{k+1} \quad (\because \text{補題1}) \\
      &= \vec{S}_{i+1} \quad (\because 1^\circ \text{に帰着})
    \end{align}
\end{enumerate}

以上 $1^\circ, 2^\circ$ から, $n=i+1$ でも ③は成立する. ゆえに③は示された. \hfill 終

\newpage

\end{document}
